\begin{PSbox}[unbreakable]{obj}{Learning Objectives}

After completing this module, students will be able to:

\begin{enumerate}
\def\labelenumi{\arabic{enumi}.}
\tightlist
\item
  Derive the distribution of Z = X + Y using convolution
\item
  Find distributions of differences, products, ratios, max, and min
\item
  Apply the bivariate Jacobian transformation method
\item
  Connect these results to engineering applications (signal+noise, reliability)
\end{enumerate}

\end{PSbox}

\PSsection{7.1}{Introduction}

Engineering systems combine random variables:

\begin{itemize}
\tightlist
\item
  \textbf{Received signal} = transmitted + noise: R = S + N
\item
  \textbf{Total interference} = sum of multiple sources: I = I\textsubscript{1} + I\textsubscript{2} + … + I\textsubscript{n}
\item
  \textbf{System lifetime} = minimum of component lifetimes: T = min(T\textsubscript{1}, T\textsubscript{2})
\item
  \textbf{SNR} = signal power / noise power: ratio of RVs
\end{itemize}

We need methods to find distributions of \textbf{functions of two (or more) RVs}.

\PSsection{7.2}{Sum of Two Random Variables: Z = X + Y}

\PSsubsection{CDF Method}

F\textsubscript{Z}(z) = P(X + Y \ensuremath{\le} z) = \ensuremath{\int}\ensuremath{\int}\emph{\{x+y\ensuremath{\le}z\} f}\{X,Y\}(x,y) dx dy

\PSsubsection{Convolution Formula (Independent Case)}

If X and Y are \textbf{independent}:

\PSdisplay{f_Z(z) = \int_{-\infty}^{\infty} f_X(x) \cdot f_Y(z-x) \, dx = (f_X * f_Y)(z)}

This is the \textbf{convolution} of f\textsubscript{X} and f\textsubscript{Y}.

\PSsubsection{Derivation}

F\textsubscript{Z}(z) = P(X + Y \ensuremath{\le} z) = \ensuremath{\int}\emph{\{-\ensuremath{\infty}\}\textsuperscript{\ensuremath{\infty}} \ensuremath{\int}}\{-\ensuremath{\infty}\}\textsuperscript{z-x} f\textsubscript{X}(x)f\textsubscript{Y}(y) dy dx\PSbr{}= \ensuremath{\int}\textsubscript{-\ensuremath{\infty}}\textsuperscript{\ensuremath{\infty}} f\textsubscript{X}(x) F\textsubscript{Y}(z-x) dx

Differentiating: f\textsubscript{Z}(z) = \ensuremath{\int}\textsubscript{-\ensuremath{\infty}}\textsuperscript{\ensuremath{\infty}} f\textsubscript{X}(x) f\textsubscript{Y}(z-x) dx \PSyes{}

\PSsubsection{Important Special Cases}

\textbf{Sum of two Gaussians (independent):}\PSbr{}X \textasciitilde{} N(\ensuremath{\mu}\textsubscript{1}, \ensuremath{\sigma}\textsubscript{1}\textsuperscript{2}), Y \textasciitilde{} N(\ensuremath{\mu}\textsubscript{2}, \ensuremath{\sigma}\textsubscript{2}\textsuperscript{2}) \ensuremath{\to} Z = X+Y \textasciitilde{} N(\ensuremath{\mu}\textsubscript{1}+\ensuremath{\mu}\textsubscript{2}, \ensuremath{\sigma}\textsubscript{1}\textsuperscript{2}+\ensuremath{\sigma}\textsubscript{2}\textsuperscript{2})

\textbf{Sum of two exponentials (same rate):}\PSbr{}X, Y \textasciitilde{} Exp(\ensuremath{\lambda}) independent \ensuremath{\to} Z \textasciitilde{} Gamma(2, 1/\ensuremath{\lambda}) = Erlang(2, \ensuremath{\lambda})

\textbf{Sum of two Poissons:}\PSbr{}X \textasciitilde{} Poisson(\ensuremath{\lambda}\textsubscript{1}), Y \textasciitilde{} Poisson(\ensuremath{\lambda}\textsubscript{2}) independent \ensuremath{\to} Z \textasciitilde{} Poisson(\ensuremath{\lambda}\textsubscript{1}+\ensuremath{\lambda}\textsubscript{2})

\textbf{Sum of two uniforms:}\PSbr{}X, Y \textasciitilde{} Uniform(0,1) independent \ensuremath{\to} Z has triangular distribution on [0,2]

\begin{PSbox}[unbreakable]{ex}{Engineering Example: Signal + Noise}

Transmitted signal S \textasciitilde{} N(A, \ensuremath{\sigma}\textsubscript{s}\textsuperscript{2}), noise N \textasciitilde{} N(0, \ensuremath{\sigma}\textsubscript{n}\textsuperscript{2}), independent.\PSbr{}Received: R = S + N \textasciitilde{} N(A, \ensuremath{\sigma}\textsubscript{s}\textsuperscript{2} + \ensuremath{\sigma}\textsubscript{n}\textsuperscript{2})

The received signal is still Gaussian — with increased variance (reduced SNR).

\end{PSbox}

\PSsection{7.3}{Difference: Z = X - Y}

For independent X, Y:\PSdisplay{f_Z(z) = \int_{-\infty}^{\infty} f_X(x) \cdot f_Y(x-z) \, dx}

\textbf{Gaussian case:} X \textasciitilde{} N(\ensuremath{\mu}\textsubscript{1}, \ensuremath{\sigma}\textsubscript{1}\textsuperscript{2}), Y \textasciitilde{} N(\ensuremath{\mu}\textsubscript{2}, \ensuremath{\sigma}\textsubscript{2}\textsuperscript{2}) \ensuremath{\to} Z = X-Y \textasciitilde{} N(\ensuremath{\mu}\textsubscript{1}-\ensuremath{\mu}\textsubscript{2}, \ensuremath{\sigma}\textsubscript{1}\textsuperscript{2}+\ensuremath{\sigma}\textsubscript{2}\textsuperscript{2})

Note: variances still ADD (not subtract) because Var(-Y) = Var(Y).

\begin{PSbox}[unbreakable]{ex}{Engineering Example: Differential Measurement}

Two sensors measure same quantity: X\textsubscript{1} = \ensuremath{\theta} + N\textsubscript{1}, X\textsubscript{2} = \ensuremath{\theta} + N\textsubscript{2}.\PSbr{}Difference: X\textsubscript{1} - X\textsubscript{2} = N\textsubscript{1} - N\textsubscript{2} \textasciitilde{} N(0, \ensuremath{\sigma}\textsubscript{1}\textsuperscript{2} + \ensuremath{\sigma}\textsubscript{2}\textsuperscript{2})

The difference eliminates the common signal but doubles the noise variance.

\end{PSbox}

\PSsection{7.4}{Product: Z = XY}

For independent X, Y:\PSdisplay{f_Z(z) = \int_{-\infty}^{\infty} \frac{1}{|x|} f_X(x) \cdot f_Y(z/x) \, dx}

\begin{PSbox}[unbreakable]{ex}{Engineering Example: Power}

Voltage V and current I in a linear circuit: P = V\ensuremath{\cdot}I.\PSbr{}If V and I have known joint distribution, we can find the power distribution.

\end{PSbox}

\PSsection{7.5}{Ratio: Z = X/Y}

For independent X, Y:\PSdisplay{f_Z(z) = \int_{-\infty}^{\infty} |y| \cdot f_X(zy) \cdot f_Y(y) \, dy}

\begin{PSbox}[unbreakable]{ex}{Engineering Example: SNR}

SNR = Signal Power / Noise Power = P\textsubscript{s} / P\textsubscript{n}

If both are chi-squared distributed (sum of squared Gaussians), the ratio follows an \textbf{F-distribution}.

\end{PSbox}

\PSsection{7.6}{Maximum: Z = max(X, Y)}

\PSsubsection{CDF Approach}

F\textsubscript{Z}(z) = P(max(X,Y) \ensuremath{\le} z) = P(X \ensuremath{\le} z AND Y \ensuremath{\le} z)

If independent: F\textsubscript{Z}(z) = F\textsubscript{X}(z) \ensuremath{\cdot} F\textsubscript{Y}(z)

PDF: f\textsubscript{Z}(z) = f\textsubscript{X}(z)F\textsubscript{Y}(z) + F\textsubscript{X}(z)f\textsubscript{Y}(z)

\begin{PSbox}[unbreakable]{ex}{Engineering Example: Parallel Redundancy}

System works if AT LEAST ONE component works. System fails only when ALL fail.\PSbr{}System lifetime = max(T\textsubscript{1}, T\textsubscript{2}) for parallel redundant components.

For T\textsubscript{1}, T\textsubscript{2} \textasciitilde{} Exp(\ensuremath{\lambda}) independent:\PSbr{}F\textsubscript{Z}(z) = (1 - e\textsuperscript{-\ensuremath{\lambda}z})\textsuperscript{2} for z \ensuremath{\ge} 0

f\textsubscript{Z}(z) = 2\ensuremath{\lambda}e\textsuperscript{-\ensuremath{\lambda}z}(1 - e\textsuperscript{-\ensuremath{\lambda}z})

Mean lifetime: E[max] = 3/(2\ensuremath{\lambda}) \textgreater{} 1/\ensuremath{\lambda} = E[single component] (reliability improved!)

\end{PSbox}

\PSsection{7.7}{Minimum: Z = min(X, Y)}

\PSsubsection{Survival Function Approach}

P(Z \textgreater{} z) = P(min(X,Y) \textgreater{} z) = P(X \textgreater{} z AND Y \textgreater{} z)

If independent: P(Z \textgreater{} z) = P(X \textgreater{} z) \ensuremath{\cdot} P(Y \textgreater{} z) = [1-F\textsubscript{X}(z)][1-F\textsubscript{Y}(z)]

CDF: F\textsubscript{Z}(z) = 1 - [1-F\textsubscript{X}(z)][1-F\textsubscript{Y}(z)]

\begin{PSbox}[unbreakable]{ex}{Engineering Example: Series System}

System fails when FIRST component fails. System lifetime = min(T\textsubscript{1}, T\textsubscript{2}).

For T\textsubscript{1} \textasciitilde{} Exp(\ensuremath{\lambda}\textsubscript{1}), T\textsubscript{2} \textasciitilde{} Exp(\ensuremath{\lambda}\textsubscript{2}) independent:\PSbr{}P(Z \textgreater{} z) = e\textsuperscript{-\ensuremath{\lambda}\textsubscript{1}z} \ensuremath{\cdot} e\textsuperscript{-\ensuremath{\lambda}\textsubscript{2}z} = e\textsuperscript{-(\ensuremath{\lambda}\textsubscript{1}+\ensuremath{\lambda}\textsubscript{2})z}

Therefore: min(T\textsubscript{1}, T\textsubscript{2}) \textasciitilde{} Exp(\ensuremath{\lambda}\textsubscript{1} + \ensuremath{\lambda}\textsubscript{2})

\textbf{Key result:} Failure rates ADD in series systems. Mean lifetime = 1/(\ensuremath{\lambda}\textsubscript{1}+\ensuremath{\lambda}\textsubscript{2}) \textless{} min(1/\ensuremath{\lambda}\textsubscript{1}, 1/\ensuremath{\lambda}\textsubscript{2}).

\end{PSbox}

\PSsection{7.8}{General Bivariate Transformation (Jacobian Method)}

\PSsubsection{Setup}

Given (X, Y) with known joint PDF, find the joint PDF of (U, V) where:

\begin{itemize}
\tightlist
\item
  U = g\textsubscript{1}(X, Y)
\item
  V = g\textsubscript{2}(X, Y)
\end{itemize}

\PSsubsection{Procedure}

\begin{enumerate}
\def\labelenumi{\arabic{enumi}.}
\tightlist
\item
  Solve for the inverse: X = h\textsubscript{1}(U, V), Y = h\textsubscript{2}(U, V)
\item
  Compute the Jacobian determinant:
\end{enumerate}

\PSdisplay{J = \begin{vmatrix} \frac{\partial x}{\partial u} & \frac{\partial x}{\partial v} \\ \frac{\partial y}{\partial u} & \frac{\partial y}{\partial v} \end{vmatrix}}

\begin{enumerate}
\def\labelenumi{\arabic{enumi}.}
\setcounter{enumi}{2}
\item
  Apply:\PSdisplay{f_{U,V}(u,v) = f_{X,Y}(h_1(u,v), h_2(u,v)) \cdot |J|}
\item
  If only U is needed, marginalize out V.
\end{enumerate}

\begin{PSbox}[unbreakable]{ex}{Example: Sum and Difference}

U = X + Y, V = X - Y \ensuremath{\to} X = (U+V)/2, Y = (U-V)/2

J = \textbar{}\ensuremath{\partial}x/\ensuremath{\partial}u \ensuremath{\cdot} \ensuremath{\partial}y/\ensuremath{\partial}v - \ensuremath{\partial}x/\ensuremath{\partial}v \ensuremath{\cdot} \ensuremath{\partial}y/\ensuremath{\partial}u\textbar{} = \textbar{}(1/2)(-1/2) - (1/2)(1/2)\textbar{} = 1/2

f\textsubscript{U,V}(u,v) = (1/2) \ensuremath{\cdot} f\textsubscript{X,Y}((u+v)/2, (u-v)/2)

\end{PSbox}

\PSsection{7.9}{MATLAB Examples}

\begin{PSbox}[unbreakable]{ex}{Example 1: Convolution — Sum of Two Uniforms}

\tcbset{PScodemode/.style={unbreakable}}

\begin{Shaded}
\begin{Highlighting}[]
\CommentTok{\%\% Z = X + Y, both Uniform(0,1)}
\VariableTok{N} \OperatorTok{=} \FloatTok{200000}\OperatorTok{;}
\VariableTok{X} \OperatorTok{=} \VariableTok{rand}\NormalTok{(}\FloatTok{1}\OperatorTok{,} \VariableTok{N}\NormalTok{)}\OperatorTok{;} \VariableTok{Y} \OperatorTok{=} \VariableTok{rand}\NormalTok{(}\FloatTok{1}\OperatorTok{,} \VariableTok{N}\NormalTok{)}\OperatorTok{;}
\VariableTok{Z} \OperatorTok{=} \VariableTok{X} \OperatorTok{+} \VariableTok{Y}\OperatorTok{;}

\VariableTok{figure}\OperatorTok{;}
\VariableTok{histogram}\NormalTok{(}\VariableTok{Z}\OperatorTok{,} \FloatTok{100}\OperatorTok{,} \SpecialStringTok{\textquotesingle{}Normalization\textquotesingle{}}\OperatorTok{,} \SpecialStringTok{\textquotesingle{}pdf\textquotesingle{}}\NormalTok{)}\OperatorTok{;}
\VariableTok{hold} \VariableTok{on}\OperatorTok{;}
\VariableTok{z} \OperatorTok{=} \VariableTok{linspace}\NormalTok{(}\FloatTok{0}\OperatorTok{,} \FloatTok{2}\OperatorTok{,} \FloatTok{200}\NormalTok{)}\OperatorTok{;}
\VariableTok{f\_theory} \OperatorTok{=} \VariableTok{max}\NormalTok{(}\FloatTok{0}\OperatorTok{,} \FloatTok{1} \OperatorTok{{-}} \VariableTok{abs}\NormalTok{(}\VariableTok{z} \OperatorTok{{-}} \FloatTok{1}\NormalTok{))}\OperatorTok{;}  \CommentTok{\% Triangular PDF}
\VariableTok{plot}\NormalTok{(}\VariableTok{z}\OperatorTok{,} \VariableTok{f\_theory}\OperatorTok{,} \SpecialStringTok{\textquotesingle{}r\textquotesingle{}}\OperatorTok{,} \SpecialStringTok{\textquotesingle{}LineWidth\textquotesingle{}}\OperatorTok{,} \FloatTok{2}\NormalTok{)}\OperatorTok{;}
\VariableTok{xlabel}\NormalTok{(}\SpecialStringTok{\textquotesingle{}Z = X + Y\textquotesingle{}}\NormalTok{)}\OperatorTok{;} \VariableTok{ylabel}\NormalTok{(}\SpecialStringTok{\textquotesingle{}PDF\textquotesingle{}}\NormalTok{)}\OperatorTok{;}
\VariableTok{title}\NormalTok{(}\SpecialStringTok{\textquotesingle{}Sum of Two Uniforms → Triangular\textquotesingle{}}\NormalTok{)}\OperatorTok{;}
\VariableTok{legend}\NormalTok{(}\SpecialStringTok{\textquotesingle{}Simulation\textquotesingle{}}\OperatorTok{,} \SpecialStringTok{\textquotesingle{}Theory\textquotesingle{}}\NormalTok{)}\OperatorTok{;}
\end{Highlighting}
\end{Shaded}

\end{PSbox}

\begin{PSbox}[unbreakable]{ex}{Example 2: Sum of Gaussians (Signal + Noise)}

\tcbset{PScodemode/.style={unbreakable}}

\begin{Shaded}
\begin{Highlighting}[]
\CommentTok{\%\% Received = Signal + Noise}
\VariableTok{mu\_s} \OperatorTok{=} \FloatTok{3}\OperatorTok{;} \VariableTok{sigma\_s} \OperatorTok{=} \FloatTok{0.5}\OperatorTok{;}  \CommentTok{\% Signal}
\VariableTok{sigma\_n} \OperatorTok{=} \FloatTok{1}\OperatorTok{;}               \CommentTok{\% Noise (zero{-}mean)}
\VariableTok{N} \OperatorTok{=} \FloatTok{100000}\OperatorTok{;}

\VariableTok{S} \OperatorTok{=} \VariableTok{mu\_s} \OperatorTok{+} \VariableTok{sigma\_s}\OperatorTok{*}\VariableTok{randn}\NormalTok{(}\FloatTok{1}\OperatorTok{,}\VariableTok{N}\NormalTok{)}\OperatorTok{;}
\VariableTok{Noise} \OperatorTok{=} \VariableTok{sigma\_n}\OperatorTok{*}\VariableTok{randn}\NormalTok{(}\FloatTok{1}\OperatorTok{,}\VariableTok{N}\NormalTok{)}\OperatorTok{;}
\VariableTok{R} \OperatorTok{=} \VariableTok{S} \OperatorTok{+} \VariableTok{Noise}\OperatorTok{;}

\VariableTok{fprintf}\NormalTok{(}\SpecialStringTok{\textquotesingle{}E[R]: Theory=\%.2f, Sim=\%.2f\textbackslash{}n\textquotesingle{}}\OperatorTok{,} \VariableTok{mu\_s}\OperatorTok{,} \VariableTok{mean}\NormalTok{(}\VariableTok{R}\NormalTok{))}\OperatorTok{;}
\VariableTok{fprintf}\NormalTok{(}\SpecialStringTok{\textquotesingle{}Var(R): Theory=\%.2f, Sim=\%.2f\textbackslash{}n\textquotesingle{}}\OperatorTok{,} \VariableTok{sigma\_s}\OperatorTok{\^{}}\FloatTok{2}\OperatorTok{+}\VariableTok{sigma\_n}\OperatorTok{\^{}}\FloatTok{2}\OperatorTok{,} \VariableTok{var}\NormalTok{(}\VariableTok{R}\NormalTok{))}\OperatorTok{;}
\VariableTok{fprintf}\NormalTok{(}\SpecialStringTok{\textquotesingle{}SNR = \%.2f dB\textbackslash{}n\textquotesingle{}}\OperatorTok{,} \FloatTok{10}\OperatorTok{*}\VariableTok{log10}\NormalTok{(}\VariableTok{mu\_s}\OperatorTok{\^{}}\FloatTok{2}\OperatorTok{/}\NormalTok{(}\VariableTok{sigma\_s}\OperatorTok{\^{}}\FloatTok{2}\OperatorTok{+}\VariableTok{sigma\_n}\OperatorTok{\^{}}\FloatTok{2}\NormalTok{)))}\OperatorTok{;}
\end{Highlighting}
\end{Shaded}

\end{PSbox}

\begin{PSbox}[unbreakable]{ex}{Example 3: Series System Reliability (Minimum)}

\tcbset{PScodemode/.style={unbreakable}}

\begin{Shaded}
\begin{Highlighting}[]
\CommentTok{\%\% Series System: T\_sys = min(T1, T2)}
\VariableTok{lambda1} \OperatorTok{=} \FloatTok{0.001}\OperatorTok{;} \VariableTok{lambda2} \OperatorTok{=} \FloatTok{0.002}\OperatorTok{;}  \CommentTok{\% Failure rates per hour}
\VariableTok{N} \OperatorTok{=} \FloatTok{100000}\OperatorTok{;}

\VariableTok{T1} \OperatorTok{=} \VariableTok{exprnd}\NormalTok{(}\FloatTok{1}\OperatorTok{/}\VariableTok{lambda1}\OperatorTok{,} \FloatTok{1}\OperatorTok{,} \VariableTok{N}\NormalTok{)}\OperatorTok{;}
\VariableTok{T2} \OperatorTok{=} \VariableTok{exprnd}\NormalTok{(}\FloatTok{1}\OperatorTok{/}\VariableTok{lambda2}\OperatorTok{,} \FloatTok{1}\OperatorTok{,} \VariableTok{N}\NormalTok{)}\OperatorTok{;}
\VariableTok{T\_sys} \OperatorTok{=} \VariableTok{min}\NormalTok{(}\VariableTok{T1}\OperatorTok{,} \VariableTok{T2}\NormalTok{)}\OperatorTok{;}

\VariableTok{fprintf}\NormalTok{(}\SpecialStringTok{\textquotesingle{}MTTF component 1: \%.0f hours\textbackslash{}n\textquotesingle{}}\OperatorTok{,} \FloatTok{1}\OperatorTok{/}\VariableTok{lambda1}\NormalTok{)}\OperatorTok{;}
\VariableTok{fprintf}\NormalTok{(}\SpecialStringTok{\textquotesingle{}MTTF component 2: \%.0f hours\textbackslash{}n\textquotesingle{}}\OperatorTok{,} \FloatTok{1}\OperatorTok{/}\VariableTok{lambda2}\NormalTok{)}\OperatorTok{;}
\VariableTok{fprintf}\NormalTok{(}\SpecialStringTok{\textquotesingle{}MTTF system (theory): \%.0f hours\textbackslash{}n\textquotesingle{}}\OperatorTok{,} \FloatTok{1}\OperatorTok{/}\NormalTok{(}\VariableTok{lambda1}\OperatorTok{+}\VariableTok{lambda2}\NormalTok{))}\OperatorTok{;}
\VariableTok{fprintf}\NormalTok{(}\SpecialStringTok{\textquotesingle{}MTTF system (sim): \%.0f hours\textbackslash{}n\textquotesingle{}}\OperatorTok{,} \VariableTok{mean}\NormalTok{(}\VariableTok{T\_sys}\NormalTok{))}\OperatorTok{;}
\end{Highlighting}
\end{Shaded}

\end{PSbox}

\begin{PSbox}[unbreakable]{ex}{Example 4: Parallel System (Maximum)}

\tcbset{PScodemode/.style={unbreakable}}

\begin{Shaded}
\begin{Highlighting}[]
\CommentTok{\%\% Parallel System: T\_sys = max(T1, T2)}
\VariableTok{lambda} \OperatorTok{=} \FloatTok{0.01}\OperatorTok{;} \VariableTok{N} \OperatorTok{=} \FloatTok{100000}\OperatorTok{;}
\VariableTok{T1} \OperatorTok{=} \VariableTok{exprnd}\NormalTok{(}\FloatTok{1}\OperatorTok{/}\VariableTok{lambda}\OperatorTok{,} \FloatTok{1}\OperatorTok{,} \VariableTok{N}\NormalTok{)}\OperatorTok{;}
\VariableTok{T2} \OperatorTok{=} \VariableTok{exprnd}\NormalTok{(}\FloatTok{1}\OperatorTok{/}\VariableTok{lambda}\OperatorTok{,} \FloatTok{1}\OperatorTok{,} \VariableTok{N}\NormalTok{)}\OperatorTok{;}
\VariableTok{T\_parallel} \OperatorTok{=} \VariableTok{max}\NormalTok{(}\VariableTok{T1}\OperatorTok{,} \VariableTok{T2}\NormalTok{)}\OperatorTok{;}

\VariableTok{fprintf}\NormalTok{(}\SpecialStringTok{\textquotesingle{}MTTF single: \%.0f hours\textbackslash{}n\textquotesingle{}}\OperatorTok{,} \FloatTok{1}\OperatorTok{/}\VariableTok{lambda}\NormalTok{)}\OperatorTok{;}
\VariableTok{fprintf}\NormalTok{(}\SpecialStringTok{\textquotesingle{}MTTF parallel (theory): \%.0f hours\textbackslash{}n\textquotesingle{}}\OperatorTok{,} \FloatTok{3}\OperatorTok{/}\NormalTok{(}\FloatTok{2}\OperatorTok{*}\VariableTok{lambda}\NormalTok{))}\OperatorTok{;}
\VariableTok{fprintf}\NormalTok{(}\SpecialStringTok{\textquotesingle{}MTTF parallel (sim): \%.0f hours\textbackslash{}n\textquotesingle{}}\OperatorTok{,} \VariableTok{mean}\NormalTok{(}\VariableTok{T\_parallel}\NormalTok{))}\OperatorTok{;}
\end{Highlighting}
\end{Shaded}

\end{PSbox}

\PSsection{7.10}{Practice Problems}

\begin{PSbox}[unbreakable]{prob}{Problem 1}

X \textasciitilde{} Exp(1), Y \textasciitilde{} Exp(1), independent. Find the PDF of Z = X + Y.

\PSsolution{} Convolution: f\textsubscript{Z}(z) = \ensuremath{\int}\textsubscript{0}\textsuperscript{z} e\textsuperscript{-x}\ensuremath{\cdot}e\textsuperscript{-(z-x)} dx = \ensuremath{\int}\textsubscript{0}\textsuperscript{z} e\textsuperscript{-z} dx = ze\textsuperscript{-z}, z\ensuremath{\ge}0.\PSbr{}This is Gamma(2,1) = Erlang(2,1). E[Z] = 2, Var(Z) = 2.

\end{PSbox}

\begin{PSbox}[unbreakable]{prob}{Problem 2}

Three components in series with failure rates \ensuremath{\lambda}\textsubscript{1}=0.001, \ensuremath{\lambda}\textsubscript{2}=0.002, \ensuremath{\lambda}\textsubscript{3}=0.003 per hour.

\begin{enumerate}
\def\labelenumi{(\alph{enumi})}
\tightlist
\item
  Find system MTTF. (b) Find P(system survives 100 hours).
\end{enumerate}

\PSsolution{} 

\begin{enumerate}
\def\labelenumi{(\alph{enumi})}
\tightlist
\item
  MTTF = 1/(0.001+0.002+0.003) = 1/0.006 = 166.7 hours
\item
  P(T\textsubscript{sys} \textgreater{} 100) = e\textsuperscript{-0.006\ensuremath{\times}100} = e\textsuperscript{-0.6} = 0.549
\end{enumerate}

\end{PSbox}

\begin{PSbox}[unbreakable]{prob}{Problem 3}

X \textasciitilde{} N(5, 4), Y \textasciitilde{} N(3, 9), independent. Find distribution and P(X + Y \textgreater{} 12).

\PSsolution{} Z = X+Y \textasciitilde{} N(8, 13). P(Z\textgreater{}12) = P(Z’ \textgreater{} (12-8)/\ensuremath{\surd}13) = Q(1.109) = 0.1337.

\emph{Next Module: {Module 8 — Conditional Expectation and Variance}}

\end{PSbox}
